% Partie : counterfactuals
% Chapitre : introduction
% Section : counterfactuals

\documentclass[../../../include/open-logic-section]{subfiles}

\begin{document}

\olfileid[fr]{cnt}{int}{cnt}

\olsection{Les conditionnels contrefactuels}

En anglais, une forme très courante et importante des constructions
« if \dots, then \dots » emploie une forme passée du subjonctif du
verbe \emph{to be} : « if it were the case that \dots, then it would
be the case that \dots ». Comme l'antécédent d'un tel conditionnel est
généralement faux, c'est-à-dire contraire aux faits, on parle de
\emph{conditionnels contrefactuels}. On les appelle aussi, en référence
à leur forme grammaticale anglaise, \emph{conditionnels au
  subjonctif}. Ils se distinguent des conditionnels \emph{à l'indicatif}, qui prennent la
forme anglaise « if it is the case that \dots, then it is the case
that \dots ». Leurs conditions de vérité diffèrent. Considérons le
célèbre exemple d'Adams :
\begin{quote}
  Si Oswald n'a pas tué Kennedy, quelqu'un d'autre l'a fait.

  Si Oswald n'avait pas tué Kennedy, quelqu'un d'autre l'aurait fait.
\end{quote}
Le premier énoncé est à l'indicatif, le second est contrefactuel. Le
premier est manifestement vrai : nous savons que le président
John~F. Kennedy a été tué par \emph{quelqu'un} et, si ce n'est pas
Lee Harvey Oswald qui l'a tué, contrairement à ce qu'affirme le
rapport Warren, c'est quelqu'un d'autre. Le second énoncé affirme
autre chose : si Oswald n'avait pas tué Kennedy, autrement dit si les
coups de feu à Dallas avaient été évités ou n'avaient pas atteint leur
but, l'histoire aurait ensuite suivi un cours tel qu'un autre
assassinat aurait réussi. Pour que cet énoncé soit vrai, il faudrait
que des forces puissantes aient conspiré pour assurer la mort de
Kennedy, comme le pensent de nombreux partisans des théories du
complot à son sujet.

La question de savoir si le conditionnel matériel rend correctement
compte du conditionnel \emph{à l'indicatif} reste débattue. Il faut
notamment déterminer si les paradoxes du conditionnel matériel
peuvent recevoir une explication compatible avec l'idée qu'il donne
les conditions de vérité des conditionnels anglais à l'indicatif. En
revanche, il est incontesté que le conditionnel matériel ne permet
pas de symboliser correctement les conditionnels contrefactuels. En
effet, bien que leurs antécédents soient généralement faux, tous les
conditionnels contrefactuels à antécédent faux ne sont pas vrais : si
l'on accepte le rapport Warren et que l'on suppose qu'il n'existait
aucun complot visant à assassiner Kennedy, le conditionnel
contrefactuel d'Adams en fournit un exemple.

Les conditionnels contrefactuels jouent un rôle important dans le
raisonnement causal : ils servent notamment à exprimer des relations
de cause à effet. Par exemple, frotter une allumette peut la faire
s'enflammer, ce que l'on peut exprimer ainsi : « Si l'on frottait
cette allumette, elle s'enflammerait. » Ni le conditionnel matériel
ni, en général, les conditionnels à l'indicatif ne permettent
d'exprimer cela : « cette allumette est frottée $\lif$ elle s'enflamme » est vrai
si l'allumette n'est jamais frottée, indépendamment de ce qui se
produirait si on la frottait. Pis encore, « cette allumette est frottée
$\lif$ elle se transforme en bouquet de fleurs » est également vrai si
elle n'est jamais frottée, alors qu'elle ne se transformerait certainement pas en
bouquet de fleurs si on la frottait.

La logique exacte des conditionnels contrefactuels reste elle aussi
débattue. Stalnaker et Lewis en ont proposé des analyses influentes.
Dans une première approximation commune à leurs analyses, on évalue
un contrefactuel « s'il se produisait~$S$, alors il se
produirait~$T$ » en examinant si~$T$ est vrai dans les situations
contrefactuelles (les « mondes possibles ») pertinentes où~$S$ est
vrai, compte tenu de leur proximité avec le monde actuel. Cette
analyse est dite « ontique », car elle fait appel à une ontologie des
mondes possibles. D'autres
analyses utilisent les probabilités conditionnelles ou des théories
de la révision des croyances. Les logiques proposées pour les
conditionnels contrefactuels sont nombreuses. Il n'existe même pas
une logique unique de Lewis et Stalnaker : bien que tous deux aient
eu recours à des mondes possibles proches du monde actuel, leurs
hypothèses distinctes conduisent à des inférences valides
différentes. Dans la sémantique des sphères étudiée plus loin,
plusieurs mondes peuvent se trouver au même degré de proximité.
Il se peut même que les sphères où l'antécédent est vrai en un
monde n'aient pas de plus petite sphère parmi elles.

\end{document}
